\documentclass[11pt]{article}

% -------------------------------------------------
% Packages
% -------------------------------------------------
\usepackage[utf8]{inputenc}
\usepackage[T1]{fontenc}
\usepackage{amsmath, amssymb, mathtools}
\usepackage{geometry}
\usepackage{booktabs}
\usepackage{hyperref}
\usepackage{enumitem}

\geometry{margin=1in}

% -------------------------------------------------
% Document
% -------------------------------------------------
\begin{document}

\begin{center}
    {\LARGE \textbf{Exercise 2.1 — Likelihood of Markov-Chain Histories}}\\[0.5em]
    {\large Recursive Macroeconomic Theory}\\[0.25em]
    {\normalsize Lars Ljungqvist and Thomas J. Sargent}
\end{center}

\vspace{1em}

\section*{Problem}

Consider the two-state Markov chain characterized by the transition matrix

\[
P =
\begin{pmatrix}
0.9 & 0.1 \\
0.3 & 0.7
\end{pmatrix},
\]
with initial distribution

\[
\pi_0 =
\begin{pmatrix}
0.5 \\
0.5
\end{pmatrix}.
\]
The observable variable \(y_t\) can take the values

\[
\bar y =
\begin{pmatrix}
1 \\
5
\end{pmatrix}.
\]
We compute the likelihood of the following observed histories for
\(t=0,1,\ldots,4\):

\begin{enumerate}[label=\alph*.]
    \item \(1,5,1,5,1\)
    \item \(1,1,1,1,1\)
    \item \(5,5,5,5,5\)
\end{enumerate}

\section*{Relevant Theory}

For a finite-state Markov chain with state history

\[
x_0,x_1,\ldots,x_T,
\]
the probability of observing the complete path is

\[
\Pr(x_0,x_1,\ldots,x_T)
=
\pi_0(x_0)
\prod_{t=0}^{T-1}
P_{x_t,x_{t+1}}.
\]
Thus, the likelihood of a particular history is obtained by multiplying

\begin{itemize}
    \item the probability of the initial state,
    \item and the transition probabilities associated with each movement along the path.
\end{itemize}
In this exercise, the observable values uniquely identify the underlying states:

\[
y_t=1 \Longleftrightarrow x_t=0,
\]
and

\[
y_t=5 \Longleftrightarrow x_t=1.
\]

\section*{Analytical Solution}

\subsection*{Part (a)}

The observed history is

\[
(1,5,1,5,1),
\]

which corresponds to the state sequence

\[
(0,1,0,1,0).
\]
Its likelihood is therefore

\[
\begin{aligned}
L_a
&=
\Pr(X_0=1)
P_{01}
P_{10}
P_{01}
P_{10} \\
&=
0.5
\times 0.1
\times 0.3
\times 0.1
\times 0.3.
\end{aligned}
\]
Hence,

\[
\boxed{
L_a = 0.00045
}
\]
This history is relatively unlikely because it requires the chain to switch states at every period. In particular, the transition from state 0 to state 1 has probability only \(0.1\).

\subsection*{Part (b)}

The observed history is

\[
(1,1,1,1,1),
\]
which corresponds to

\[
(0,0,0,0,0).
\]
The likelihood is

\[
\begin{aligned}
L_b
&=
\Pr(X_0=1)
P_{00}^4 \\
&=
0.5(0.9)^4.
\end{aligned}
\]
Therefore,

\[
\boxed{
L_b = 0.32805
}.
\]
This path is relatively likely because state 0 is highly persistent:

\[
P_{00}=0.9.
\]

\subsection*{Part (c)}

The observed history is

\[
(5,5,5,5,5),
\]
corresponding to the state sequence

\[
(1,1,1,1,1).
\]
Its likelihood is

\[
\begin{aligned}
L_c
&=
\Pr(X_0=1)
P_{11}^4 \\
&=
0.5(0.7)^4.
\end{aligned}
\]
Hence,

\[
\boxed{
L_c = 0.12005
}.
\]
State 1 is also persistent, but less so than state 0 because

\[
P_{11}=0.7 < P_{00}=0.9.
\]
Therefore,

\[
\boxed{
L_b > L_c > L_a
}.
\]
As we can see, the ranking reflects the transition matrix's persistence structure; indeed, remaining in state 0 is the most likely of the three histories, remaining in state 1 is less likely, and repeated switching between the two states is substantially less likely.

\section*{Log-Likelihood Representation}

For longer histories, multiplying many probabilities can lead to numerical underflow. It is therefore often convenient to work with the log-likelihood.
In more detail, the path likelihood

\[
L
=
\pi_0(x_0)
\prod_{t=0}^{T-1}
P_{x_t,x_{t+1}}
\]
can be rewritten as

\[
\log L
=
\log \pi_0(x_0)
+
\sum_{t=0}^{T-1}
\log P_{x_t,x_{t+1}}.
\]
The original likelihood can then be recovered through

\[
L = \exp(\log L),
\]
which is a way more useful representation both in numerical and econometric applications.


\section*{Extension}

A natural computational extension is to simulate repeated histories from the Markov chain and compare their empirical frequencies with the theoretical path probabilities.

This provides a direct numerical illustration of how the transition matrix determines state persistence first, path probabilities right afterwards, and lastly empirical frequencies.

\section*{Reference}

Ljungqvist, L. and Sargent, T. J. (2018).
\textit{Recursive Macroeconomic Theory},
4th edition, MIT Press.

\end{document}